% =====================================================================
% chapters/04_Bench.tex — LaTeX rendering of ../sections/04_bench.en.md (2026-09-24 00:01:41)
% Section 4 of the SHORT paper. \input from main.tex, after 03_Agent.tex.
% This file DEFINES \label{sec:bench}, referenced by the reading map of Section 1.4.
% Tables: one. tab:scale is set as table* across both columns (../PLAN.md § 9: Table 1 is
%   full width), caption ABOVE, booktabs (loaded by aamas.cls). Fifteen rows and five
%   columns do not fit \columnwidth.
% Figures: none.
% Display equation: one, unnumbered (\[ ... \]), amsmath being loaded by aamas.cls. The three
%   other metrics live in Appendix C, in the supplementary material.
% PLACEHOLDER, DELIBERATE, waiting on lot B of ticket 103: one cell reading
%   "[pending]" in the last column of Table 1, as the master writes it. It is set in \textbf{} since
%   23 September 2026, at the author's request: everything that does not yet exist is bold,
%   so that it is found at a glance in the compiled PDF. It is NOT \todo. Do not remove
%   it before the replay lands.
% The row "Expert prompt, typed classifier (in sample)" carries its reservation in the row
%   label and nowhere else, by the review (§ 9 bis). It enters no ranking.
% Citations: tisseo2023emc2 for the survey (twice in 4.1), horl2021eqasim, lin1991divergence,
%   rubner2000earth, ke2017lightgbm, zhu2005kernel, breiman2001random, as in the master.
% Appendices A, B, C and D are named in plain text: they live in the supplementary material,
%   outside this project, so there is no \label to \ref.
% Derived file: regenerate after every validated pass of ../sections/04_bench.en.md.
% HAND-CARRIED on 23 September 2026, three passages, the generator md_vers_tex.py refusing
%   this chapter (its Markdown carries a float caption, which the MD model counts as a
%   paragraph: 23 blocks against the 21 of this file). The three are the eqasim sentence of
%   4.1, the availability paragraph of 4.1 and the fifth criterion of 4.4.
% AUTHOR-ADDED FOOTNOTE — PRESERVE ON REGENERATION: the PROGEDO archive URL in 4.1, set as
%   a footnote on the model of the survey footnote of the abstract. It is deliberately NOT
%   in ../sections/04_bench.en.md, where it would cost body words in an over-budget section.
% Citation added on 23 September 2026: horl2021eqasim, the synthesis chain that produced the
%   pool the cohort is drawn from. The key was already in sample.bib and cited nowhere.
% REALIGNED on 23 September 2026, 4.4: the four criteria are named, and "situational
%   principles" becomes "general criteria", the single term the paper uses (R12). This file
%   still carried the wording the author corrected on 22 September.
% Author's edit of 23 September 2026 on the master, carried here: the archive is named
%   "PROGEDO-Diffusion", not "Quetelet-PROGEDO-Diffusion".
% Author's edits of 23 September 2026, second batch, carried here:
%   - 4.3 names the five strata, split three and two so the enumeration stays inside R1,
%     and says which two are ordered. The gloss on "stratum" goes with it.
%   - 4.4 loses "The text was tuned against one language model on a separate calibration
%     population, then served unchanged to the others (Appendix D)." That sentence rested on
%     the contradiction between the two masters of 21 and 22 September, flagged in the
%     section report. With it goes the only mention of Appendix D in this file.
% HAND-CARRIED again on 24 September 2026, whole chapter, against the master of
%   2026-09-24 00:01:41: the opening and closing sentences and the caption's last sentence,
%   absent from the master, are gone; "Section 3.2" is \ref{sec:receives}, "Table 1" is
%   \ref{tab:scale}.
% HAND-CARRIED on 24 September 2026, three passages and nothing else, so the date on line 2
%   stays where it was: this file is still behind its master.
%   - 4.1: "in an anonymised repository", with its address as a footnote. The address comes
%     from ../repository.tex (\repobase, \repodisplay), the only place where its TBD is
%     written. PRESERVE ON REGENERATION, like the PROGEDO footnote.
%   - 4.4: "That prompt adds a fifth general criterion, real exposure to the weather." was
%     wrong. The classifier's prompt has no weather bullet; it rewrites the chain-friction
%     bullet on two sides (Appendix D.4 prints it). Replaced by "That prompt rewrites the
%     criterion on walking and waiting, so that a long direct walk also counts as a cost."
%   - 4.4: "Appendix D gives the three prompts in full, with one trip decided under each."
%   NOT carried: the master's rewrite of the same 4.4 paragraph around a "calibration cohort"
%   (author's decision of 24 September), which its own source comment makes true only once
%   the classifier's prompt is re-tuned on that cohort.
% =====================================================================

\section{The comparison bench}
\label{sec:bench}

\subsection{The cohort}
\label{sec:cohort}

Every score in this paper comes from one sealed cohort of 1{,}000 synthetic personas. A persona
is one synthetic person with a household and a day of trips. They form 499 whole households,
inside the study area. The eqasim synthesis chain \citep{horl2021eqasim}
produced the pool of 11{,}329 people they were drawn from. It builds that pool from census and
national travel survey data. Their reference remains the local survey \citep{tisseo2023emc2},
in which residents describe every trip of the day before.

Thirteen margins of the cohort are controlled against the survey. All thirteen conform. Each
margin is the share of one trait, such as age class or car ownership. The equivalence bound
was one point, set before measuring. The largest gap is 0.50 point (Appendix A).

The cohort travels slightly less than the survey population: 3.30 trips per persona and 3.69
per persona who travels at all, against 3.53 and 3.95. That makes 3{,}299 trips per day.

The weather and the transport offer are those of a weekday of the survey's collection period,
from September 2022 to February 2023. That weekday is drawn for each persona, so the cohort
covers the whole period.

The survey records one mode per trip, not an itinerary. Every score therefore sums the
probability of the options that share a mode.

We release the protocol, the code, the seeds and the sealed cohort in an anonymised
repository\footnote{\href{\repobase}{\texttt{\repodisplay}}}. Our research convention
forbids passing on the survey microdata themselves. The survey is obtained from the
PROGEDO-ADISP archive, on individual
request\footnote{\url{https://www.progedo.fr/donnees-disponibles/donnees-localisees/}}
\citep{tisseo2023emc2}.

\subsection{The protocol in three rules}
\label{sec:rules}

Three rules close three ways of biasing a comparison between a generative agent and a
tabular model.

Rule 1 gives every decision-maker the same 21 variables, twelve for the person and
household, three for the trip, six for geometry (Appendix B).

Rule 2 scores the probability mass a decision-maker puts on each option, not the option it
ranks first. The option ranked first is not the decision the simulation plays
(Section~\ref{sec:receives}).

Rule 3 restricts every tabular prediction to the modes actually offered for that trip, then
rescales it to 100\,\%. A tabular model predicts over every mode, whereas a language model sees
only the options offered. The share of a mode is split evenly among its options.

We equalise the 21 variables, not the information each side holds. The language-based
decision-makers receive three things no tabular model can receive. They see the agenda of
the remaining trips, the weather of the coming hours, and each itinerary step by step. The
tabular references received something the language-based decision-makers never did, the
39{,}203 survey trips on which they were estimated. The exposure asymmetry runs the other way,
which makes those models high references.

\subsection{What we measure}
\label{sec:quantities}

We measure at two scales, the modal split of the cohort and the agreement on each trip. The
modal split is the share of trips made by car, public transport, walking and
cycling. The survey gives the reference share globally, then inside five strata. Three are
nominal, occupation, gender and trip purpose. Two are ordered, age class and distance class.

One composite number carries the aggregate reading. It adds the global divergence to the
weighted mean divergence of each stratum.
\[
  \mathcal{C}_{\text{EMD--JSD}} = \mathrm{JSD}^{\text{global}}
  + \sum_{d\ \text{nominal}} w_d\,\overline{\mathrm{JSD}}^{\,d}
  + \sum_{d\ \text{ordinal}} w_d\,\overline{\mathrm{EMD}}^{\,d}
\]
The divergence is Jensen--Shannon \citep{lin1991divergence} in base 2 on nominal strata.
Ordered strata take the earth mover's distance \citep{rubner2000earth}, which respects the
class order. Both lie in $[0, 1]$, the distance once divided by the number of class steps, and
both are multiplied by 100. A nominal stratum averages its classes, each weighted by its number
of personas. An ordered one averages the four modes, each weighted by its survey share. The
weights are set by design, not fitted. The global term weighs 1, age, occupation and purpose
0.5 each, gender and distance 0.3 each. A lower composite means a distribution closer to the
survey.

Two other readings accompany the composite, all three published together. The first
recomputes it without the trips that offered a single option, where no preference is
expressed. The second is the L1 error on the global shares, the sum of absolute gaps in
percentage points.

A difference between two decision-makers counts only when it exceeds what a new cohort would
move. Another cohort of 1{,}000 personas built the same way would shift a composite by
$\pm$1.3 points. Resolution varies by decision-maker, from 0.9 point for the typed classifier
under its expert prompt to 2.1 under its minimal prompt.

Trips of one person are not independent, so every interval comes from resampling clusters at
the person level. A paired difference compares two decision-makers on the persons both scored,
at 95\,\%.

The unit-level audit asks the other question, on the days respondents actually described. It
covers 9{,}621 trips of 2{,}930 respondents, each carrying its declared mode. We publish
accuracy, cross-entropy in nats, and recall and precision per mode. Cross-entropy falls as a
decision-maker puts more probability on the declared mode. Recall is the share of a mode's
declared trips that it designates.

\subsection{Floors, references, decision-makers}
\label{sec:deciders}

Three floors give the level a decision-maker reaches without behavioural knowledge. One
draws uniformly over the options offered, another puts everything on the car, the area's
majority mode. The third, the minimum-duration heuristic, keeps the fastest option offered.

Four tabular references give the level that estimation on the survey reaches. They are a
multinomial logit, the standard discrete-choice model, and three learned models. These are
LightGBM gradient boosting \citep{ke2017lightgbm}, a kernel logistic regression
\citep{zhu2005kernel} and a random forest \citep{breiman2001random}. We estimate the four at
strict parity, on the same survey partition. None of them leads on all three readings, so we
read the ceiling one reading at a time.

Three families of decision-makers are under test. The minimal prompt gives a language model
the facts of the trip and the output format, with no general criteria. The expert prompt
keeps that format and draws the model's attention to four general criteria. Two bear on
comfort, the walking and waiting public transport really adds, and an unhurried pace for an
older traveller. The other two bear on constraints, heavy shopping bags to carry and a
working day with no slack. The third family is a zero-shot classifier with typed output
(TypeSafe System One, \texttt{jev-1.13.0}). It reads that same text and returns one
probability per option without writing a sentence.

We obtained the expert prompt by successive mutations of one text, under three constraints.
Each iteration reads the gaps per stratum of the prompt in service, proposes one targeted
rewrite, and keeps it when the over-represented mode recedes. The constraints forbid a
numeric threshold, forbid naming a survey category, and require scoring on verbalised
distributions. We call tuned agent the expert prompt served to \texttt{gemini-3.5}, the
best-scoring language-model condition in Table~\ref{tab:scale}.

The classifier has its own expert prompt, tuned on the first cohort. That prompt rewrites
the criterion on walking and waiting, so that a long direct walk also counts as a cost.
Every crossing of models and prompts is therefore measured on a second cohort, which shares
no persona with the first. Appendix D gives the three prompts in full, with one trip decided
under each.

\begin{table*}[tp]
  \caption{The fifteen decision-makers on the three readings, with the inter-seed range where it was measured.}
  \label{tab:scale}
  \begin{tabular}{lrrrl}
    \toprule
    Decision-maker & Composite & Excluding single option & L1 on global shares & Inter-seed range \\
    \midrule
    Uniform random                                     & 50.16 & 57.90 & 86.80 & deterministic \\
    All-car                                            & 30.73 & 28.41 & 58.00 & deterministic \\
    Minimum duration                                   & 26.97 & 23.93 & 53.62 & deterministic \\
    Minimal prompt, \texttt{mistral-large}             & 14.75 & 20.72 & 44.71 & not replayed \\
    Minimal prompt, typed classifier                   & 13.75 & 19.84 & 42.76 & not replayed \\
    Minimal prompt, \texttt{gemini-3.1}                & 12.38 & 16.65 & 39.88 & not replayed \\
    Minimal prompt, \texttt{gemini-3.5}                &  7.02 & 10.39 & 24.08 & \textbf{[pending]} \\
    Expert prompt, \texttt{gemini-3.1}                 &  8.98 & 12.39 & 31.25 & not replayed \\
    Expert prompt, \texttt{mistral-large}              &  7.63 & 12.27 & 26.64 & not replayed \\
    Expert prompt, \texttt{gemini-3.5}                 &  4.86 &  6.86 & 13.85 & 0.56 \\
    Multinomial logit                                  &  4.02 &  6.63 &  8.99 & deterministic \\
    Random forest                                      &  4.09 &  5.81 &  5.28 & deterministic \\
    Expert prompt, typed classifier (in sample)        &  3.65 &  6.58 & 10.48 & 0.05 \\
    Kernel logistic regression                         &  3.61 &  5.61 &  6.69 & deterministic \\
    Gradient boosting                                  &  3.60 &  5.84 &  9.49 & deterministic \\
    \bottomrule
  \end{tabular}
\end{table*}
